\input{template.tex}

\ifPDFTeX
  \renewcommand{\rmdefault}{cmr}
\fi

\def\SheetNo{5}
\setlength{\emergencystretch}{2em}

\begin{document}
\sheettitle{Комбинаторика.}

\begin{problems}

\item Сколькими способами можно распределить $15$ различных гостей по четырём различным комнатам, заняв каждую комнату хотя бы одним гостем?

\item Имеется $r$ различных предметов и $n+p$ человек. Докажите, что число распределений предметов, при которых каждый из заранее указанных $n$ человек получает хотя бы один предмет, равно
\[
\sum_{j=0}^{n}(-1)^j C_n^j(n+p-j)^r.
\]

\item Сколькими способами можно заполнить таблицу $k\times l$ числами $+1$ и $-1$, чтобы произведение в каждой строке и в каждом столбце было равно $1$?

\item Каково наибольшее возможное число различных точек пересечения диагоналей $n$-угольника?

\item На выборах один кандидат получил $12$ голосов, а другой --- $8$. Сколько существует порядков подсчёта голосов, при которых после каждого подсчитанного голоса первый кандидат строго опережает второго? Порядки различаются последовательностью голосов за двух кандидатов.

\item Пусть $E_n$ --- число перестановок чисел $1,\ldots,n$, в которых ровно один раз следующее число больше предыдущего. Например, такова перестановка $4312$. Докажите, что при каждом натуральном $n$
\[
E_{n+1}=2E_n+n.
\]

\item Сколько целых чисел в промежутке от $1$ до $999$ включительно не делятся ни на $3$, ни на $5$, ни на $7$?

\item Пусть $p$ --- простое число и $0<k<p$. Докажите, что $p\mid C_p^k$:
\begin{parts}
  \item с помощью формулы для биномиального коэффициента;
  \item комбинаторно, без использования формул.
\end{parts}

\item На полке стоят $n$ книг. Сколькими способами можно выбрать $p$ из них, если между соседними выбранными книгами и после последней выбранной должно оставаться хотя бы $s$ книг? Считайте, что $n\ge(s+1)p$.

\item Найдите коэффициент при $x^{10}y^{10}z^{10}$ в многочлене
\[
(x+y)^{10}(y+z)^{10}(z+x)^{10}.
\]

\item Докажите, что каждое из следующих количеств равно $n$-му числу Каталана $c_n$:
\begin{parts}
  \item число заполнений таблицы $2\times n$ числами $1,\ldots,2n$, каждое из которых используется один раз, с возрастанием слева направо в строках и сверху вниз в столбцах;
  \item число способов соединить $2n$ точек окружности попарно $n$ хордами так, чтобы хорды не пересекались;
  \item число параллеломино периметра $2n+2$. Параллеломино задаётся парой ломаных по линиям клетчатой сетки: они начинаются в $(0,0)$, имеют общий конец, идут только вверх и вправо и не имеют других общих точек. Периметр --- сумма длин двух ломаных.
\end{parts}

\item Докажите, что число корневых упорядоченных полных бинарных деревьев с $2n+1$ вершиной равно числу Каталана $c_n$. Здесь у каждой вершины либо нет потомков, либо есть ровно два --- левый и правый; порядок потомков учитывается. Постройте явную биекцию с правильными скобочными последовательностями из $n$ пар скобок.

\item На плоскости выбраны $n$ точек: никакие три не лежат на одной прямой, а никакие четыре --- на одной окружности. Проведены все прямые через пары выбранных точек и все окружности через тройки. Каково наибольшее возможное число различных точек пересечения этих прямых с этими окружностями?

\item Во всех клетках доски $2020\times2020$ расставлены буквы В, С, О и Ш. Расстановка называется гармоничной, если в каждом квадрате $2\times2$ все четыре буквы различны. Найдите число гармоничных расстановок.

\item Имеются $12$ карандашей попарно различной длины. Сколькими способами можно уложить их в коробку в два слоя по шесть карандашей так, чтобы в каждом слое карандаши были упорядочены по возрастанию длины слева направо, а каждый карандаш верхнего слоя лежал строго над карандашом нижнего слоя и был короче него?

\item Троих англичан, троих немцев и троих русских ставят в ряд. Сколько существует порядков, в которых никакие три соотечественника не стоят подряд? Все девять человек различны.

\item Сколько чисел от $1$ до $1\,000\,000$ включительно не являются ни квадратами, ни кубами, ни пятыми степенями целых чисел?

\item На прямой отмечены $100$ точек. Сколькими способами можно выбрать из них $10$ точек так, чтобы никакие две выбранные точки не были соседними среди отмеченных?

\item Докажите, что число разбиений натурального числа $n$ на не более чем $m$ натуральных слагаемых равно числу его разбиений на натуральные слагаемые, каждое из которых не превосходит $m$. Порядок слагаемых не учитывается.

\item Дан клетчатый квадрат $101\times101$. Внутри него выбирается квадрат $100\times100$, внутри выбранного --- квадрат $99\times99$, и так далее до квадрата $1\times1$. Оказалось, что последний квадрат совпадает с центральной клеткой исходного квадрата. Сколько существует таких последовательностей квадратов? Ответ не должен содержать знака многоточия.

\item На полоске из $n$ клеток фишка стоит в крайней левой клетке. За ход её можно передвинуть на одну или две клетки вправо. Докажите, что число способов добраться до крайней правой клетки равно $F_n$, где $F_1=F_2=1$ --- первые числа Фибоначчи.

\item Сколькими способами можно раскрасить клетки доски $20\times20$ в три цвета так, чтобы все цвета были использованы и в каждом уголке из трёх клеток присутствовали ровно два цвета?

\item Фишка находится в левой нижней клетке доски $m\times n$. За ход её можно передвинуть на одну клетку вверх или вправо. Сколько существует путей в правую верхнюю клетку?

\item Сколько раз встречается каждая из цифр $0,1,\ldots,9$ при записи всех натуральных чисел от $1$ до $10^n-1$ без ведущих нулей?

\item Инверсией в перестановке называется пара чисел, в которой большее стоит раньше меньшего. Найдите суммарное число инверсий во всех перестановках чисел $1,2,\ldots,n$.

\end{problems}
\end{document}
